%% KOMA

\usepackage[
  paper=a4paper,
  left=25mm,
  right=25mm,
  top=25mm,
  bottom=20mm,
  includefoot,
  foot=\baselineskip,
  bindingoffset=0mm
]{geometry}

\usepackage[utf8]{inputenc} % encoding
\usepackage[T1]{fontenc} % selecting font encodings
\usepackage[ngerman]{babel} % German hyphenation
\usepackage[babel,german=quotes]{csquotes} % German quotes
\usepackage{lmodern} % high resolution font
\renewcommand*\familydefault{\sfdefault} % sans-serif font
% text companion fonts (for symbols like µ, bullet, musical note, etc.)
\usepackage{textcomp}

\usepackage{scrlayer-scrpage}
\clearpairofpagestyles
\setkomafont{pageheadfoot}{\sffamily\footnotesize}
\setkomafont{pagehead}{\bfseries}
\setkomafont{pagination}{}

\KOMAoptions{
  headsepline = true,
  footsepline = false,
  plainfootsepline = false
}
\renewcommand*{\sectionmarkformat}{}

\automark{section}
\ihead{\myName}
\chead{\myDocumentType}
\ohead{\headmark}
\ofoot*{\pagemark}

\usepackage{scrdate} % date handling
\usepackage{scrhack} % included because \float@listhead is not
% supported anymore
% http://www.bakoma-tex.com/doc/latex/koma-script/scrhack.pdf page 3

%% other packages

\usepackage{etoolbox}

\usepackage[onehalfspacing]{setspace}

% different table style (TODO: try!)
% \usepackage{booktabs}
\usepackage{multirow} % combining rows/columns in tables
\usepackage{makecell}
% \usepackage{tabularx}
\usepackage{longtable} % allows tables that continue to the next page

\usepackage{paralist} % for \begin{compactitem}

% TODO: probably use minted instead, might not be necessary then?
\usepackage{listings}
\lstset{
  basicstyle=\footnotesize, % Global Code Style
  numbers=left,
  numberstyle=\tiny,
  columns=flexible,
  breaklines=true,
  postbreak=\mbox{\textcolor{red}{$\hookrightarrow$}\space},
}
% from
% https://tex.stackexchange.com/questions/303465/higher-asterisks-in-lstlisting-environment
\makeatletter
\lst@CCPutMacro
\lst@ProcessOther {"2A}{\raisebox{2pt}{*}}
\@empty\z@\@empty
\makeatother

% modified version of IEEEtran with support for online sources (url and date)
% => register as @misc and add url = {...} and optionally urldate = {...}
% \bibliographystyle{IEEEtran_online}

\usepackage[style=ieee,backend=biber]{biblatex}

% better \includegraphics and float objects
\usepackage{graphicx}
\usepackage{float}
% \usepackage[abs]{overpic} % adding content on top of images

\usepackage{siunitx}
% \sisetup{scientific-notation = true}
\sisetup{exponent-product = \ensuremath{\cdot}}
\sisetup{product-symbol = \ensuremath{\cdot}}
\sisetup{detect-all}
\sisetup{per-mode = symbol}
\sisetup{locale=DE}
\usepackage{eurosym}
\DeclareSIUnit{\sieuro}{\mbox{\euro}}

% \usepackage[colorinlistoftodos,prependcaption,textsize=footnotesize]{todonotes}

\usepackage{adjustbox} % interface to transform latex content
\usepackage{subcaption}

% \usepackage{lscape} % placing certain parts of a document in landscape
% \usepackage{tikz} % creating graphics in latex source using a
% drawing interface
% \usepackage{pgfplots} % generating plots in latex source using tikz
% \pgfplotsset{compat=1.16}
% \usepackage{pgfplotstable} % loading and displaying tables of numerical data

%% glossary

\usepackage[acronym,nonumberlist,nopostdot]{glossaries}
\usepackage{glossary-mcols}
\setglossarystyle{mcolindex}
\glsenablehyper % enable hyperlinks to glossary in text
\renewcommand{\glossarypreamble}{\glsfindwidesttoplevelname[\currentglossary]}
% funktioniert nicht mit mcolindex
\makenoidxglossaries

\usepackage[
  hidelinks,              % Das Links nicht farbig
  bookmarks,              % Bookmarks erstellen
  bookmarksopen,          % Bookmarks bei öffnen des PDF anzeigen
  bookmarksnumbered,      % damit Nummern in Bookmarks
  bookmarksopenlevel = 1, % damit nur die erste Ebene der Bookmarks
  % beim öffnen angezeigt werden
  pdftoolbar = false,
  pdfmenubar = false,
  pdftitle={    \myDocumentType~über~\myTitle~von~\myName},
  pdfsubject={},
  pdfkeywords={},
  pdfauthor = {    \myName},
  pdfstartpage = {  1}
]{hyperref}

% set names of references for \autoref (i.e. subsection => "Abschnitt")
\addto\extrasngerman{\def\subsectionautorefname{Abschnitt}}
\addto\extrasngerman{\def\subsubsectionautorefname{Abschnitt}}
\addto\extrasngerman{\def\appendixautorefname{Anhang}}
\addto\extrasngerman{\def\lstlistingautorefname{Codeausschnitt}}
\renewcommand{\lstlistingname}{Codeausschnitt}

%% appendix

\usepackage{appendix}
\renewcommand{\appendixname}{Anhang}
\renewcommand{\appendixtocname}{Anhang}
\renewcommand{\appendixpagename}{Anhang}
\newcommand{\Appendixautorefname}{Anhang}
\newcommand*\appendixmore{
  \clearpage
  %\addsec{{\appendixname}}%
  \renewcommand{\thesubsection}{{\Alph{subsection}}}
}

\renewcommand{\k}[1]{\textit{#1}}
\newcommand{\f}[1]{\textbf{#1}}

\let\oldautoref\autoref
\let\oldref\ref
\renewcommand{\autoref}[1]{\oldautoref{#1}}
\renewcommand{\ref}[1]{\oldref{#1}}

% inserts the reference and the name (e.g. "Abschnitt 1 Einleitung")
\newcommand{\fullref}[1]{\autoref{#1} \nameref{#1}}

%% miscellaneous

% add equation number to the current line (e.g. in "alignat" blocks)
\newcommand\eqnumberthis{\addtocounter{equation}{1}\tag{\theequation}}
% define a command alias for a math expression (usage:
% \newmath{\<name>}{<math>})
\newcommand{\newmath}[2]{\newcommand{#1}{\ensuremath{#2}}}
